%% =====================================================================
%%  第2章 準備
%%  書くこと：本研究の前提となる既存の知識（扱う問題の定義，用いる手法の基礎）
%%  ここは自分の提案ではない．教科書や論文の内容を，出典を明記して簡潔にまとめる．
%%  自分の提案は「提案手法」の章に書き，既存の内容とはっきり分ける．
%%  以下では例として，0-1 ナップサック問題とその貪欲法を説明している．
%%  自分の研究に合わせて書き換える．
%% =====================================================================
\chapter{準備}
\label{chap:preliminaries}

本章では，本研究で扱う問題を定義し，提案手法の基礎となる既存の解法について述べる．
\begin{writingnote}
この章には，提案手法と数値実験を理解するために必要な既存の知識だけを書く．教科書のように網羅する必要はない．
定義や既存の解法は，原典（元の論文や標準的な教科書）を引用して説明する．
文献によって定義が異なる用語は，本論文で用いる定義を明記し，以後はその意味だけで使う．
記号は分野の標準的な教科書に合わせ，一度定義した記号を別の意味で使わない．主な記号は\cref{chap:notation}のように一覧にするとよい．
\end{writingnote}

\section{対象とする問題}
\label{sec:problem}

%% 問題は，言葉による説明 → 数式による定式化 → 記号の説明 の順に書くとわかりやすい．
%% 複数行の数式は align 環境で書き，& を置いた位置で揃える．&& の後には条件を書く．
%% 連続した式は \crefrange{最初}{最後} で「式(2.1)〜式(2.3)」のように参照できる．
例として，0-1 ナップサック問題~\cite{kellerer2004knapsack}を考える．
0-1 ナップサック問題は，価値と重さが与えられた $n$ 個の品物の中から，
重さの合計が容量を超えないように品物を選び，価値の合計を最大化する問題であり，
\crefrange{eq:knapsack-objective}{eq:knapsack-binary}のように定式化できる．
\begin{align}
  \text{maximize}   \quad & \sum_{i=1}^{n} v_i x_i \label{eq:knapsack-objective} \\
  \text{subject to} \quad & \sum_{i=1}^{n} w_i x_i \le W \label{eq:knapsack-capacity} \\
                          & x_i \in \{0, 1\} && i = 1, \dots, n \label{eq:knapsack-binary}
\end{align}
ここで，$v_i$ と $w_i$ はそれぞれ品物 $i$ の価値と重さ，$W$ は容量であり，
$x_i$ は品物 $i$ を選ぶとき 1，選ばないとき 0 をとる決定変数である．
\cref{eq:knapsack-objective}は選んだ品物の価値の合計を最大化することを，
\cref{eq:knapsack-capacity}は重さの合計が容量を超えないことを表す．
数理最適化の基礎については，文献~\cite{kanatani2005optimization}などを参照されたい．

%% 記号が少ない場合は，このように式の直後に「ここで，……」と本文中で定義する．
%% 記号が多い場合（大きな数理モデルなど）は，定式化の前に記号を「集合」「パラメータ」
%% 「決定変数」に分けて，description 環境で一覧にするとよい．例：
%%   \begin{description}
%%     \item[集合] $\mathcal{I} = \{1, \dots, n\}$ は品物の集合である．
%%     \item[パラメータ] $v_i$ と $w_i$ は品物 $i$ の価値と重さ，$W$ は容量である．
%%     \item[決定変数] $x_i$ は品物 $i$ を選ぶとき 1，選ばないとき 0 をとる．
%%   \end{description}

\section{既存の解法}
\label{sec:greedy}

\begin{writingnote}
既存の解法の考え方と手順を，その解法を知らない読者にもわかるように説明し，長所と短所を述べる．
短所は，後の章で述べる課題や提案手法の工夫につながるように具体的に書く．
既存の解法に自分で加えた変更は，この章ではなく「提案手法」の章に書く．
\end{writingnote}

%% アルゴリズム（擬似コード）は algorithm 環境で書く（algorithm2e パッケージ）．
%% 各行の最後には \; を付ける．
0-1 ナップサック問題の代表的な近似解法に，価値と重さの比 $v_i / w_i$ が大きい順に品物を選ぶ貪欲法がある~\cite{dantzig1957discrete}．
貪欲法の手順を\cref{alg:greedy}に示す．
貪欲法は計算が速い一方で，最適解が得られるとは限らない．

\begin{algorithm}[tbp]
  \caption{0-1 ナップサック問題に対する貪欲法}
  \label{alg:greedy}
  \KwIn{品物の価値 $v_1, \dots, v_n$，重さ $w_1, \dots, w_n$，容量 $W$}
  \KwOut{選んだ品物の集合 $S$}
  品物を $v_i / w_i$ の大きい順に並べ替える\;
  $S \gets \emptyset$，$c \gets 0$ \tcp*{$c$ は選んだ品物の重さの合計}
  \ForEach{並べ替えた順の品物 $i$}{
    \If{$c + w_i \le W$}{
      $S \gets S \cup \{i\}$\;
      $c \gets c + w_i$\;
    }
  }
  \KwRet{$S$}\;
\end{algorithm}
